% ============================================================================
% pandoc-preamble.tex
% Paquetes y macros que requiere el LaTeX generado por Pandoc a partir de los
% archivos Markdown del informe (tablas booktabs/longtable, bloques de código
% con resaltado de sintaxis, listas compactas, etc.).
%
% La clase MIA-USA.cls ya carga: graphicx, xcolor[table], longtable, array,
% colortbl, tabularx, amsmath, tikz, fancyhdr, multirow, multicol, tcolorbox.
% Aquí solo se añade lo que falta. Se usa pdfLaTeX (no se importan fuentes
% Times New Roman/fontspec del metadata.yaml, que eran para XeLaTeX).
% ============================================================================

% --- Tablas generadas por Pandoc -------------------------------------------
\usepackage{booktabs}        % \toprule \midrule \bottomrule
\usepackage{calc}            % cálculos de ancho de columnas de Pandoc
\providecommand{\real}[1]{#1}

% Pandoc define columnas de tablas con tipos personalizados; estos comandos
% reproducen su plantilla por defecto para que las tablas compilen.
\newlength{\cslhangindent}
\setlength{\cslhangindent}{1.5em}

% Pandoc envuelve las tablas sin título en {\def\LTcaptype{none} ...}. Como la
% plantilla carga el paquete caption (vía subcaption), este intenta usar un
% contador 'none' que no existe. Lo definimos para evitar el error fatal.
\makeatletter
\@ifundefined{c@none}{\newcounter{none}}{}
\makeatother

% --- Tablas anchas: que quepan en la página --------------------------------
\usepackage{etoolbox}
\AtBeginEnvironment{longtable}{\footnotesize\setlength{\tabcolsep}{3pt}\renewcommand{\arraystretch}{0.95}}

% Tablas en orientación horizontal cuando el ancho lo requiere
\usepackage{pdflscape}

% --- Bloques de código con resaltado (salida de Pandoc) --------------------
\usepackage{fvextra}
% Ajuste automático de línea para que el código no se desborde del margen
\DefineVerbatimEnvironment{Highlighting}{Verbatim}{breaklines,breakanywhere,commandchars=\\\{\}}
\newcommand{\VerbBar}{|}
\newcommand{\VERB}{\Verb[commandchars=\\\{\}]}
\newenvironment{Shaded}{}{}
\newcommand{\AlertTok}[1]{\textcolor[rgb]{1.00,0.00,0.00}{\textbf{#1}}}
\newcommand{\AnnotationTok}[1]{\textcolor[rgb]{0.38,0.63,0.69}{\textbf{\textit{#1}}}}
\newcommand{\AttributeTok}[1]{\textcolor[rgb]{0.49,0.56,0.16}{#1}}
\newcommand{\BaseNTok}[1]{\textcolor[rgb]{0.25,0.63,0.44}{#1}}
\newcommand{\BuiltInTok}[1]{\textcolor[rgb]{0.00,0.50,0.00}{#1}}
\newcommand{\CharTok}[1]{\textcolor[rgb]{0.25,0.44,0.63}{#1}}
\newcommand{\CommentTok}[1]{\textcolor[rgb]{0.38,0.63,0.69}{\textit{#1}}}
\newcommand{\CommentVarTok}[1]{\textcolor[rgb]{0.38,0.63,0.69}{\textbf{\textit{#1}}}}
\newcommand{\ConstantTok}[1]{\textcolor[rgb]{0.53,0.00,0.00}{#1}}
\newcommand{\ControlFlowTok}[1]{\textcolor[rgb]{0.00,0.44,0.13}{\textbf{#1}}}
\newcommand{\DataTypeTok}[1]{\textcolor[rgb]{0.56,0.13,0.00}{#1}}
\newcommand{\DecValTok}[1]{\textcolor[rgb]{0.25,0.63,0.44}{#1}}
\newcommand{\DocumentationTok}[1]{\textcolor[rgb]{0.73,0.13,0.13}{\textit{#1}}}
\newcommand{\ErrorTok}[1]{\textcolor[rgb]{1.00,0.00,0.00}{\textbf{#1}}}
\newcommand{\ExtensionTok}[1]{#1}
\newcommand{\FloatTok}[1]{\textcolor[rgb]{0.25,0.63,0.44}{#1}}
\newcommand{\FunctionTok}[1]{\textcolor[rgb]{0.02,0.16,0.49}{#1}}
\newcommand{\ImportTok}[1]{\textcolor[rgb]{0.00,0.50,0.00}{\textbf{#1}}}
\newcommand{\InformationTok}[1]{\textcolor[rgb]{0.38,0.63,0.69}{\textbf{\textit{#1}}}}
\newcommand{\KeywordTok}[1]{\textcolor[rgb]{0.00,0.44,0.13}{\textbf{#1}}}
\newcommand{\NormalTok}[1]{#1}
\newcommand{\OperatorTok}[1]{\textcolor[rgb]{0.40,0.40,0.40}{#1}}
\newcommand{\OtherTok}[1]{\textcolor[rgb]{0.00,0.44,0.13}{#1}}
\newcommand{\PreprocessorTok}[1]{\textcolor[rgb]{0.74,0.48,0.00}{#1}}
\newcommand{\RegionMarkerTok}[1]{#1}
\newcommand{\SpecialCharTok}[1]{\textcolor[rgb]{0.25,0.44,0.63}{#1}}
\newcommand{\SpecialStringTok}[1]{\textcolor[rgb]{0.73,0.40,0.53}{#1}}
\newcommand{\StringTok}[1]{\textcolor[rgb]{0.25,0.44,0.63}{#1}}
\newcommand{\VariableTok}[1]{\textcolor[rgb]{0.10,0.09,0.49}{#1}}
\newcommand{\VerbatimStringTok}[1]{\textcolor[rgb]{0.25,0.44,0.63}{#1}}
\newcommand{\WarningTok}[1]{\textcolor[rgb]{0.38,0.63,0.69}{\textbf{\textit{#1}}}}

% --- Listas compactas y varios de Pandoc -----------------------------------
\providecommand{\tightlist}{%
  \setlength{\itemsep}{0pt}\setlength{\parskip}{0pt}}
\setlength{\emergencystretch}{3em} % evita líneas desbordadas

% Caracteres de dibujo de cajas (─ │ └ ├ ┌ ┐ …) usados en los árboles de
% directorios dentro de los bloques de código. pmboxdraw los define para
% pdfLaTeX+inputenc.
\usepackage{pmboxdraw}

% Otros caracteres Unicode que la fuente base / inputenc no traen.
\usepackage{newunicodechar}
\newunicodechar{⋮}{\ensuremath{\vdots}}        % U+22EE elipsis vertical
\newunicodechar{→}{\ensuremath{\rightarrow}}   % U+2192 flecha derecha
\newunicodechar{←}{\ensuremath{\leftarrow}}    % U+2190
\newunicodechar{↔}{\ensuremath{\leftrightarrow}}
\newunicodechar{−}{\ensuremath{-}}             % U+2212 signo menos
\newunicodechar{≤}{\ensuremath{\leq}}          % U+2264
\newunicodechar{≥}{\ensuremath{\geq}}          % U+2265
\newunicodechar{×}{\ensuremath{\times}}        % U+00D7 (por si acaso)
\newunicodechar{►}{\ensuremath{\blacktriangleright}} % U+25BA
\newunicodechar{✅}{\checkmark}                 % U+2705
\newunicodechar{❌}{\ensuremath{\times}}        % U+274C
\newunicodechar{⁴}{\textsuperscript{4}}        % U+2074 superíndices (notas)
\newunicodechar{⁵}{\textsuperscript{5}}
\newunicodechar{⁶}{\textsuperscript{6}}
\newunicodechar{⁷}{\textsuperscript{7}}
\newunicodechar{⁸}{\textsuperscript{8}}
\newunicodechar{⁹}{\textsuperscript{9}}
\usepackage{amssymb}                            % \checkmark, \blacktriangleright

% Pandoc puede emitir \passthrough en código en línea
\providecommand{\passthrough}[1]{#1}

% Enlaces de imágenes/figuras de Pandoc dentro de \includegraphics
\makeatletter
\providecommand{\pandocbounded}[1]{#1}
\makeatother
